\section{Experimental Results}
\label{sec:results}

\begin{figure*}[t] 
  \centering
  
  % ==================== ROW 1: LATENT SPACES ====================
  \begin{subfigure}{0.24\textwidth}
    \includegraphics[width=\linewidth]{csp_latent_space}
    \caption{CSP Components}
  \end{subfigure}\hfill
  \begin{subfigure}{0.24\textwidth}
    \includegraphics[width=\linewidth]{TSNE_EEGNet_Latent_Space}
    \caption{EEGNet $t$-SNE}
  \end{subfigure}\hfill
  \begin{subfigure}{0.24\textwidth}
    \includegraphics[width=\linewidth]{TSNE_Psd_Latent_Space}
    \caption{EEGPsdNet $t$-SNE}
  \end{subfigure}\hfill
  \begin{subfigure}{0.24\textwidth}
    \includegraphics[width=\linewidth]{TSNE_DualBranch_Latent_Space}
    \caption{EEGConnNet $t$-SNE}
  \end{subfigure}
  
  \vspace{1.0em} % Clear vertical separation between the two rows
  
  % ==================== ROW 2: CONFUSION MATRICES ====================
  \begin{subfigure}{0.24\textwidth}
    \includegraphics[width=\linewidth]{csp_confusion_matrix}
    \caption{CSP Baseline CM}
  \end{subfigure}\hfill
  \begin{subfigure}{0.24\textwidth}
    \includegraphics[width=\linewidth]{eeg_confusion_matrix}
    \caption{EEGNet Baseline CM}
  \end{subfigure}\hfill
  \begin{subfigure}{0.24\textwidth}
    \includegraphics[width=\linewidth]{psd_confusion_matrix}
    \caption{EEGPsdNet CM}
  \end{subfigure}\hfill
  \begin{subfigure}{0.24\textwidth}
    \includegraphics[width=\linewidth]{conn_confusion_matrix}
    \caption{EEGConnNet CM}
  \end{subfigure}

  \caption{Comprehensive comparison of pre-decision feature geometries (top row) and resulting classification error distributions (bottom row) across the evaluated frameworks. Moving from left to right illustrates how the sequential integration of deep temporal filters and multimodal priors progressively disentangles the complex intra-limb overlap observed in the classical spatial baseline.}
  \label{fig:master_grid}
\end{figure*}

\subsection{Latent Representations and Error Dynamics}
\label{sec:latent_confusion}
To understand the performance disparities between the classical pipeline and the proposed deep learning architectures, we directly compared their pre-decision latent spaces alongside their resulting confusion matrices on the withheld validation set. This comparative $4 \times 2$ grid is presented in Figure~\ref{fig:master_grid}.

The classical CSP+LDA baseline (Fig.~\ref{fig:master_grid}a, b) highlights the fundamental difficulty of intra-limb classification. CSP spatial filters for all classes (\emph{elbow}, \emph{hand}, \emph{forearm} and \emph{rest}) are heavily overlapped in the primary component space. Consequently, the linear classifier suffers from severe off-diagonal confusion, failing to establish robust boundaries between states.

Transitioning to the convolutional baseline, EEGNet (Fig.~\ref{fig:master_grid}c, d) demonstrates the capacity of temporal receptive fields to begin untangling these clusters. The $t$-SNE projection reveals distinct cluster formation, significantly improving global accuracy. However, noticeable bridging remains between classes, resulting in a misclassification bottleneck.

The integration of multimodal priors does not resolve these boundary ambiguities. While EEGConnNet (Fig.~\ref{fig:master_grid}g, h) force tighter intra-class clustering and wider inter-class margins, EEGPsdNet remains very similar to the baseline.

\subsection{Overall Performance Benchmarks}
\label{sec:accuracy}
The quantitative impact of these feature space improvements is summarized in Table~\ref{tab:main}. Models were evaluated based on maximum validation accuracy during optimization, with the classical baseline fitted purely on the training partition to ensure an unbiased comparison. 

The baseline CSP approach achieved approximately 36\% accuracy, reinforcing its limitations for same-limb paradigms. The baseline temporal CNN elevated this to $\sim$56\%. However, the dual-branch architectures incorporating late-fusion spectral embeddings (EEGPsdNet) and dynamic cross-spectral attention (EEGConnNet) achieved the highest generalization, both scoring between 61\% and 62\%. This confirms that supplementing raw voltage dynamics with structural network connectivity yields superior decoding bounds for complex motor imagery.

\begin{table*}[htbp]
  \centering
  \begingroup
  \captionsetup{font=normalsize}
  \caption{Held-out four-class classification accuracy, parameter footprint, and architectural notes evaluated on the withheld validation fraction.}
  \label{tab:main}
  \normalsize
  \setlength{\tabcolsep}{10pt}
  \renewcommand{\arraystretch}{1.55}
  \begin{tabularx}{\textwidth}{@{} >{\raggedright\arraybackslash}p{3.5cm} >{\raggedright\arraybackslash}X >{\centering\arraybackslash}p{2.2cm} >{\centering\arraybackslash}p{2.2cm} @{}}
    \toprule
    \textbf{Model Name} & \textbf{Description \& Representations} & \textbf{Model Size} & \textbf{Accuracy [\%]} \\
    \midrule
    \textbf{Block I: CSP+LDA}~\cite{gramfort2013meg} & 
    Motor-band CSP features classified via LDA. & 
    10KB & 
    \(\sim\)36 \\
    
    \textbf{Block II: EEGNet}~\cite{lawhern2018eegnet} & 
    Temporal CNN baseline operating directly on filtered, cue-locked voltage epochs using separable convolutions. & 
    15KB & 
    \(\sim\)56 \\
    
    \textbf{Block II: EEGPsdNet} & 
    Late fusion architecture combining temporal waveform embeddings with Welch PSD maps ($3$--$30\,\text{Hz}$). & 
    39KB & 
    \(\sim\)61 \\
    
    \textbf{Block II: EEGConnNet} & 
    Dual-branch processing raw waveforms and cross-spectral coherence matrices. & 
    160KB & 
    \(\sim\)61\,--\,62 \\
    \bottomrule
  \end{tabularx}
  \endgroup
\end{table*}